\chapter*{Segunda nota de actualización, 10 de septiembre de 2026}
\addcontentsline{toc}{chapter}{Segunda nota de actualización, 10 de septiembre}

\titular{El índice de la Gaceta \textbf{se reescribió por tercera vez y esta vez se
corrigió}: lo que ahora publica coincide renglón por renglón con el mapa de anexos que este
documento reconstruyó el día 9 leyendo la primera página de cada PDF. Y se descargó el archivo de las \textbf{43.636 AGEB} de la
declaratoria urbana: sus cuentas cierran, pero \textbf{sólo sobre la clave de localidad
actual a junio de 2026}. Sobre las claves que el propio archivo trae al frente, sus
\textbf{2.423 municipios y 4.531 localidades} son \textbf{2.416 y 4.540}.}

\textbf{Qué es esta nota y qué no es.} Es la segunda adenda, y obedece a la misma regla que
la primera: \textbf{el cuerpo de los dieciséis capítulos no se reescribe}. Lo que se añade
aquí no corrige ninguna cifra del documento ---no hay ninguna que corregir--- sino que
\textbf{confirma con la fuente} una reconstrucción que el día 9 era nuestra, y añade una
pieza de datos que el paquete anunciaba y no entregaba. \textbf{Lo que estaba fuera de
alcance sigue fuera de alcance}, por tercer día.

\section*{Uno. Tercer día sin Ruta A}

Los seis recursos que abren la Ruta A se sondearon por tercera vez, el 10 de septiembre a las
21:52. \textbf{Ninguno cambió}: los cuatro analíticos del proyecto responden \textbf{404} con
el mismo cuerpo de 2.978 bytes de las dos sondas anteriores, y las dos ranuras del portal de
datos abiertos \textbf{503} con el mismo cuerpo de 592 bytes. La ranura de plazas responde
200 y sigue sirviendo el archivo del ejercicio 2026, con la misma fecha de modificación y el
mismo tamaño al byte.

\textbf{La distinción que este documento hizo el día 8 se sostiene tres días después}: no es
que el árbol de los analíticos esté vacío, es que \textbf{no existe}; y no es que falte el
recurso del portal de datos, es que \textbf{el servicio está caído}. Son dos fallas distintas
y se diagnostican distinto. En consecuencia siguen sin poder hacerse el diff institucional,
la auditoría de etiquetado de los anexos transversales, la composición por programa
presupuestario, la distribución territorial, y la separación entre atención médica y
pensiones en el IMSS y el ISSSTE.

\section*{Dos. El índice de la Gaceta se corrigió solo, y confirma el cuadro \ref{cua:anexos}}

La primera nota dejó registrado que el índice del día de la entrega se había reescrito y que
las letras se habían reasignado, y publicó el mapa real de los trece anexos \textbf{leído de
la primera página de cada PDF}, porque el índice no servía para identificarlos. El 10 de
septiembre el índice se reescribió por tercera vez: \textbf{el bloque que el día 9 estaba
comentado ---con las letras del año anterior y fechado 2026--- se descomentó, se le
corrigieron las letras y se le corrigió el año}.

\begin{table}[htbp]\centering\scriptsize
\caption{Lo que el índice de la Gaceta dijo de cada anexo, en sus tres versiones}\label{cua:convergencia}
\begin{tabular}{l>{\raggedright\arraybackslash}p{3.2cm}>{\raggedright\arraybackslash}p{2.9cm}>{\raggedright\arraybackslash}p{2.9cm}>{\raggedright\arraybackslash}p{2.9cm}}
\toprule
Letra & Contenido, leído del PDF & Índice del 8 sep. & Índice del 9 sep. & Índice del 10 sep. \\
\midrule
\textbf{F} & Ley Aduanera & anunciada como Código Fiscal & Ley Aduanera & Ley Aduanera \\
\textbf{G} & Ley de Economía Digital & anunciada como informe arancelario & Ley de Economía Digital & Ley de Economía Digital \\
\textbf{H} & Ley Catastral y Registral & no listada & Ley Catastral y Registral & Ley Catastral y Registral \\
\textbf{J} & Informe arancelario, art. 131 & no listada; iba como G & Informe arancelario, art. 131 & Informe arancelario, art. 131 \\
\textbf{K} & Nota metodológica y declaratoria ZAP 2027 & no listada & comentada, como H y fechada 2026 & Nota metodológica y declaratoria ZAP \\
\textbf{L} & Listado ZAP rurales 2027 & no listada & comentada, como J y fechada 2026 & Listado ZAP rurales 2027 \\
\textbf{M} & Listado ZAP urbanas 2027 & no listada & comentada, como K y fechada 2026 & Listado ZAP urbanas 2027 \\
\textbf{N} & Estimación de subsidios de vivienda 2027 & no listada & comentada, como L y fechada 2026 & Estimación de subsidios de vivienda 2027 \\
\bottomrule
\end{tabular}
\\[2pt]\begin{minipage}{0.92\textwidth}\scriptsize Las tres copias del índice están archivadas, con marca de tiempo. Los anexos A a E no aparecen aquí porque su descripción no cambió en ninguna de las tres versiones. \textbf{La columna del día 10 coincide renglón por renglón con el cuadro \ref{cua:anexos}}, que se construyó dos días antes leyendo la primera página de cada PDF. El índice del 10 conserva un error de la fuente: titula el anexo K «zonas de atención prioritaria 2026» en su primera cláusula y 2027 en la segunda.\end{minipage}
\end{table}


\porqueimporta{\textbf{La columna del día 10 es, renglón por renglón, el cuadro
\ref{cua:anexos}}, que se construyó dos días antes y sin ella. La reconstrucción propia deja
de ser reconstrucción propia y pasa a ser \textbf{reconstrucción confirmada por la fuente}.
Pero la lección de método no se ablanda, se afila: \textbf{el índice converge, y tarda dos
días}. Un documento que se escribe el día de la entrega no puede esperarlos. \textbf{El día
de la corrida la única fuente que identifica un anexo es el PDF}, y la regla de la primera
nota ---identifica por la primera página, guarda el índice con marca de tiempo cada vez---
queda en pie con una adición: \textbf{vuelve al tercer día y compara; la cuarta copia es la
confirmación que la fuente te debía}.}

Dos detalles menores, para que el registro esté completo. El índice del 10 estrena un
\textbf{anexo O} ---el orden del día de la sesión--- y un \textbf{anexo V} ---una
comunicación de la Mesa Directiva sobre turnos---, además de una sección de prevenciones:
\textbf{nada de eso es paquete económico} y nada de eso toca el perímetro de este documento.
Y el índice conserva un error de la fuente: describe el anexo K como «zonas de atención
prioritaria 2026» en su primera cláusula y 2027 en la segunda, cuando el PDF es de 2027 en
las dos.

\section*{Tres. Las 43.636 AGEB de la declaratoria urbana}

La primera nota cerró anotando que la nota metodológica del anexo K \textbf{publica una ruta
de descarga} de las 43.636 AGEB urbanas con sus variables, y que este proyecto no la había
usado. Se descargó el 10 de septiembre: un archivo comprimido de 2,1 MB con una sola hoja de
cálculo dentro, doce columnas y 43.636 renglones de datos.

\subsection*{3.1 No es una pieza de la cola del paquete}

\textbf{El archivo se publicó el 15 de julio de 2026}, según la fecha de modificación que
declara el servidor: \textbf{dos meses antes de la entrega del paquete}. No llegó con la cola
ni con el paquete. \textbf{Lo que llegó con el paquete es el nombre del archivo}, porque el
conteo de AGEB va dentro del nombre ---cambia cada año--- y sólo la nota metodológica lo
dice.

\porqueimporta{Hay aquí una asimetría de publicación: \textbf{el dato estuvo disponible
dos meses antes que el documento que lo cita}, y era inalcanzable porque nadie
podía adivinar su dirección. No es un problema de calendario: es de catálogo. Para la
corrida del año que viene esto cambia una prioridad: \textbf{la ruta puede intentarse
antes del día de la entrega}, en cuanto se sepa el conteo, y no esperar a la cola.}

\subsection*{3.2 Las cuentas cierran, y sobre qué clave}

\begin{table}[htbp]\centering\footnotesize
\caption{Las cuentas de la declaratoria urbana contra el archivo de las 43.636 AGEB}\label{cua:zapcuentas}
\begin{tabular}{>{\raggedright\arraybackslash}p{5.0cm}>{\raggedleft\arraybackslash}p{2.2cm}>{\raggedleft\arraybackslash}p{2.0cm}>{\raggedleft\arraybackslash}p{2.2cm}}
\toprule
Cuenta & Declarado en el anexo K & Sobre la clave actual & Sobre las claves del archivo \\
\midrule
AGEB urbanas prioritarias & 43.636 & 43.636 & 43.636 \\
Entidades federativas & 32 & 32 & 32 \\
\textbf{Municipios} & 2.423 & 2.423 & \textbf{2.416} \\
\textbf{Localidades urbanas} & 4.531 & 4.531 & \textbf{4.540} \\
\midrule AGEB dentro de los 1.575 municipios ZAP rurales & no se publica & 25.836 & --- \\
\bottomrule
\end{tabular}
\\[2pt]\begin{minipage}{0.92\textwidth}\scriptsize Fuente: anexo K de la Gaceta 7121 para lo declarado; y \texttt{2027/2027\_zap-urbanas-agebs\_variables.zip}, descargado el 10 de septiembre de 2026, para lo calculado. \textbf{Los dos renglones en negritas son el aviso}: municipios y localidades sólo cierran contra la última columna del archivo, \emph{clave de localidad actual a junio de 2026}. Quien cuente sobre las columnas de clave que el archivo trae al frente publicará 2.416 y 4.540. El traslape con la lista rural sí lo publica la nota metodológica: el archivo lo reproduce al AGEB.\end{minipage}
\end{table}


Las cuatro cuentas de la declaratoria cierran contra el archivo: 43.636 áreas
geoestadísticas básicas \textbf{sin una sola repetida} y 32 entidades federativas. Los
municipios y las localidades \textbf{también cierran, pero no sobre las columnas de clave
que el archivo trae al frente}: sobre ésas dan 2.416 municipios y 4.540 localidades.
Cierran sobre la última columna, \emph{clave de localidad actual a junio de 2026}.

\textbf{La nota metodológica advierte esa columna, y hay que decirlo con precisión.} Su
párrafo final explica que la clave actualizada a junio de 2026 se agrega al anexo «con fines
operativos», porque la creación de nuevos municipios puede haber movido la clave de localidad
respecto de la cartografía de 2020, y declara que las AGEB se identifican \textbf{con su clave
oficial de origen}. \textbf{Lo que la nota no dice es cuál de las dos claves cuenta}: que sus
2.423 municipios y sus 4.531 localidades están calculados sobre la columna actualizada y no
sobre la de origen, que es la que el archivo pone al frente. Eso es lo que aquí se reconstruye.

La diferencia son \textbf{120 AGEB reasignadas}. Ciento cinco cambian de municipio y quince
cambian de localidad dentro del mismo municipio ---localidades absorbidas por la cabecera, en
Cuernavaca, Mazatlán, Tijuana, Uruapan---. De las que cambian de municipio, \textbf{102 se
van a siete municipios que la clave vieja no tiene}.

\begin{table}[htbp]\centering\footnotesize
\caption{Los siete municipios que sólo existen bajo la clave actualizada}\label{cua:zapnuevos}
\begin{tabular}{ll>{\raggedright\arraybackslash}p{2.4cm}>{\raggedright\arraybackslash}p{4.0cm}r}
\toprule
Clave & Municipio & Entidad & Se desprende de & AGEB \\
\midrule
02007 & San Felipe & Baja California & 02002 Mexicali & 16 \\
04013 & Dzitbalché & Campeche & 04001 Calkiní & 8 \\
12082 & Las Vigas & Guerrero & 12053 San Marcos & 8 \\
12085 & San Nicolás & Guerrero & 12023 Cuajinicuilapa & 4 \\
24059 & Villa de Pozos & San Luis Potosí & 24028 San Luis Potosí & 18 \\
25019 & Eldorado & Sinaloa & 25006 Culiacán & 25 \\
25020 & Juan José Ríos & Sinaloa & 25001 Ahome; 25011 Guasave & 23 \\
\midrule \textbf{Suma} & & & & \textbf{102} \\
\bottomrule
\end{tabular}
\\[2pt]\begin{minipage}{0.92\textwidth}\scriptsize Claves y AGEB, del archivo de variables; nombres, del anexo L, que marca a los siete en su columna de \emph{municipios de reciente creación}. Las 102 AGEB de este cuadro más 3 que pasan de Chinameca a Oteapan ---Veracruz, municipios que ya existían--- son las 105 que cambian de municipio; las 15 restantes de las 120 reasignadas cambian de localidad dentro del mismo municipio, absorbidas por la cabecera.\end{minipage}
\end{table}


\textbf{Los siete son municipios de creación reciente, y el anexo L lo confirma}: los marca a
los siete en su columna de \emph{municipios de reciente creación}, donde hay nueve en total
---los dos restantes, Ñuu Savi y Santa Cruz del Rincón, en Guerrero, no tienen AGEB urbanas---.
Los nueve son los que la nota metodológica llama «municipios de creación posterior al año
2020».

\textbf{Aquí hay dos sietes distintos y no deben confundirse.} La cadena acumulativa de la
nota pasa de 1.638 a 1.645 municipios rurales porque, de esos nueve, dos ya eran prioritarios
y \textbf{se añaden siete}. El de este cuadro es otro: los siete que \textbf{tienen AGEB
urbanas} y que sólo existen bajo la clave actualizada. Las dos cifras coinciden en el número y
\textbf{no se ha verificado que coincidan en los municipios}; este documento no lo afirma.

\porqueimporta{\textbf{Quien cuente sobre las columnas del frente publicará 2.416 municipios
y no sabrá por qué no coincide con la declaratoria.} El archivo trae dos mapas municipales
---el del censo y el de junio de 2026--- y la declaratoria cuenta sobre el segundo sin
decirlo. Para este proyecto la consecuencia es concreta: \textbf{cualquier cruce futuro entre
la ZAP y el gasto tiene que hacerse sobre la clave actualizada}, o siete municipios de
creación reciente ---y las 102 AGEB que viven en ellos--- quedarán colgados de un municipio
que ya no existe.}

\subsection*{3.3 El traslape entre las dos declaratorias}

\textbf{25.836 de las 43.636 AGEB urbanas ---el 59,2~\%--- caen dentro de los 1.575 municipios
de la lista ZAP rural.} La cifra \textbf{la publica la nota metodológica}, y el archivo la
reproduce al AGEB: la columna que marca el traslape parte el universo completo y da
exactamente 25.836. La nota añade el desglose que el archivo no trae: de esas 25.836,
\textbf{17.035 ya eran prioritarias} por otros criterios y \textbf{8.801 se agregaron} por
éste, hasta el universo de 43.636.

\porqueimporta{\textbf{Las dos declaratorias no son disjuntas: se traslapan en seis de cada
diez AGEB urbanas.} Sumar «municipios ZAP rurales» y «municipios con AGEB urbanas ZAP» para
obtener un universo territorial cuenta doble, y el traslape no es marginal sino mayoritario.
Este documento no cruza la ZAP con el gasto ---hacerlo exige la distribución territorial, que
sigue enganchada a la Ruta A--- pero el denominador queda ahora \textbf{en carpeta y
recalculado}, no sólo citado.}

\section*{Cuatro. Qué se volvió a correr y qué no}

Esta nota, como la anterior, \textbf{no reescribe ningún capítulo}, y las verificaciones que
validan el cuerpo siguen siendo las que se publicaron. Lo que se corrió otra vez:

\begin{itemize}
\item \textbf{Identidades: de 41 a 49 pruebas, ninguna falla.} Las ocho nuevas son las cuatro
cuentas de la declaratoria contra el archivo, la comprobación de que ninguna AGEB está
repetida, la de que la columna del traslape parte el universo completo, la del traslape
mismo ---25.836 AGEB, contadas en el archivo contra las declaradas en la nota--- y la de su
desglose, 17.035 ya prioritarias más 8.801 agregadas. El guion no lleva las cifras dentro:
viven en el bloque de restituciones del ejercicio, como manda la separación del archivo.
\item \textbf{Compuerta semántica}, aplicada al titular, a los tres pies de cuadro y a las
cuatro frases de «por qué importa» de esta nota. \textbf{Ninguna cifra modificada, y tres
frases reescritas, las tres por la misma falla y la peor de las que la compuerta persigue}:
el borrador atribuía a la fuente un silencio que no tenía. Decía que la nota metodológica no
advierte la clave actualizada ---la advierte, en su párrafo final---; que el traslape de
25.836 AGEB no está publicado ---está publicado, y con desglose---; y ligaba nuestros siete
municipios con AGEB urbanas a los siete que la nota suma en su cadena rural, que son otro
conjunto. \textbf{Las tres se corrigieron contra el texto del anexo K, no contra la memoria
del redactor.}
\item \textbf{Checklist de omisiones}: el renglón B13 ---la ZAP como denominador
territorial--- pasa de «registrada, no cruzada» a \textbf{«denominador en mano, cruce
pendiente»}. El cruce sigue enganchado a la Ruta A. \textbf{B1, la vara de suficiencia en
salud, sigue fuera por tercer ejercicio}: es una línea de texto y un cociente.
\end{itemize}

Con los tres cuadros de esta nota, los del documento son \textbf{veinticinco}.
\textbf{Ninguna cifra del cuerpo cambió, ni el 9 ni el 10 de septiembre.}

\begin{ficha}{segunda nota de actualización}
\fila{Perímetro}{Ninguno presupuestal. Esta nota no toca gasto ni ingreso: verifica un
denominador territorial y el mapa documental del paquete.}
\fila{Línea de comparación}{Ninguna. No hay variación real que calcular.}
\fila{Año de los pesos}{No aplica: no hay pesos en esta nota.}
\fila{Denominadores}{Los conteos de AGEB, municipios, localidades y entidades se citan como
los publica la fuente, \textbf{y se recalculan contra el archivo de origen}. El traslape ---25.836 AGEB, 59,2~\%--- lo publica la nota
metodológica y aquí se recalcula contra el archivo.}
\fila{Fuentes}{Archivo de variables de las ZAP urbanas, de la Secretaría de Bienestar,
descargado el 10 de septiembre de 2026 de la ruta que publica el anexo K y archivado en
la carpeta del ejercicio 2027; el nombre exacto está al pie del cuadro
\ref{cua:zapcuentas}. Anexos K y L de la Gaceta Parlamentaria 7121; y las tres copias
fechadas del índice de la Gaceta, del 8, el 9 y el 10 de septiembre.}
\fila{Qué no puede afirmarse con ellas}{Nada sobre el gasto. La ZAP es un denominador
territorial y \textbf{este documento no lo ha cruzado con un peso}: para hacerlo hace falta
la distribución territorial del gasto, que exige los analíticos del proyecto. Tampoco que
los 2.416 municipios de las columnas del archivo sean un error: son otro mapa municipal,
el del censo, y la declaratoria no cuenta sobre él.}
\end{ficha}
